%=============================================================================!
% 17.6  LITHIUM ON THE MOVE                                                   !
%=============================================================================!
\section{Lithium on the move}
\label{sec:ip-hops}

A marble in an egg box that is being shaken rattles in its cup and now
and then jumps into the next. Counting its jumps needs two decisions:
what a cup is, and how long the marble must stay in a new cup for the
visit to count, since a marble that touches the rim of a neighbour and
falls back has gone nowhere. If the egg box itself slides about on the
table, the cups move with it, and the marble's cup must be read in the
frame of the box. A lithium atom in a gallery of graphite, the gap
between two neighbouring sheets, is such a marble. This section counts its hops in dilute graphite, one lithium
among 64 carbon atoms, and turns them into rates with error bars, taking
Chapter~13's hops over a model surface (\cref{sec:th-langevin}) and
Chapter~16's counting of events (\cref{sec:ob-reactions}) to a real
material.

\subsection*{The sites}

Relaxed at \SI{0}{\kelvin}, lithium sits over the centre of a hexagon of
the sheet below, a hollow, and under an atom of the sheet above, since in
graphite's A, B stacking a hollow of one sheet lies under an atom of the
next. Graphite has an inversion centre midway between an atom of the
sheet below and the atom of the sheet above that lies over it: sending
every point to the point as far beyond that centre on the opposite side
swaps the two sheets and leaves the crystal as it was. It carries
lithium's site to one over an atom of the sheet below, a corner of
lithium's hexagon, and under a hollow of the sheet above: a second family
of sites with the same energy, $a/\sqrt3 = \SI{1.42}{\angstrom}$ from the
first, for the lattice constant $a = \SI{2.4648}{\angstrom}$ of the
relaxed dilute cell. Each family is the hollows of one sheet, a
triangular grid, and the two grids are offset by $a/\sqrt3$, so together
they form a honeycomb on which the families alternate, as the two kinds
of atom of a sheet do (\cref{ex:ip-honeycomb}).

Moved along the straight line to the nearest site of the other family,
with its position in the plane held at each of eight steps while its
height and the carbon atoms relax, lithium's energy rises to
\SI{73}{\milli\electronvolt} halfway and falls back to zero at the far
site, as the inversion requires; along the line to the next site of its
own family, across the middle of a C--C bond, it rises to
\SI{107}{\milli\electronvolt} (\cref{fig:ip-hops}a). Without D3, at the
same relaxed points, the two are 53 and \SI{94}{\milli\electronvolt}. One
carbon atom of each sheet is held in the plane as well, since otherwise
the sheets would slide to carry their sites back under lithium. Every
path between two sites must cross the ridge that divides them, whose
lowest point, the saddle, sets the barrier (\cref{sec:en-gradient}), so
the highest point of any one path is at least the barrier, and
\SI{73}{\milli\electronvolt} bounds it from above. The relaxation of the
carbon atoms is part of every site. With them held where they relaxed
around the first site, the other family's site lies
\SI{136}{\milli\electronvolt} higher and is no minimum at all, the
straight way to it passing 56, 123 and \SI{111}{\milli\electronvolt}, so
a map over a frozen host finds only half the sites.

\subsection*{Sheets that slide}

Each run, at 300, 450 and \SI{600}{\kelvin}, lasted \SI{20}{\pico\second}
under CSVR with $\tau_{\mathrm T} = \SI{1}{\pico\second}$ and steps of
\SI{1}{\femto\second}, keeping lithium's position every
\SI{10}{\femto\second} and every atom's every \SI{100}{\femto\second}.
The cell was the one relaxed at \SI{0}{\kelvin}, held fixed, which at
these temperatures holds the sheets closer than zero pressure would, as
LiC$_6$'s pressure along $c$ in its own cell showed (\cref{sec:ip-lic6});
the rates below belong to that cell. In them the sheets slide as wholes. The sheet
below lithium slides by as much as 0.96, 2.40 and \SI{1.59}{\angstrom},
and the sheet above by as much as 1.94, 4.76 and \SI{2.97}{\angstrom}
against it, comparable with or further than the \SI{1.42}{\angstrom}
between neighbouring sites, so a site fixed in the frame of the cell is
soon in the wrong place. Each family is the hollows of one sheet, the
first those of the sheet below and the second those of the sheet above.
Lithium's position is therefore read in the frame of each sheet in turn,
the sheet's slide being the mean displacement of its atoms, interpolated
between the frames that hold them, and lithium is assigned to the
nearest site of either family. Sites fixed in the cell where the sheets
started give a count of hops that falls steadily as the least stay of the
next subsection grows, from 85 to 20 at \SI{450}{\kelvin}.

How far the sheets wander is a property of the small cell. Near its best
stacking, sliding a sheet of $N$ atoms by $s$ along one direction costs
about $\frac12N\kappa s^2$, with $\kappa$ the stiffness for one atom,
since every atom of the sheet changes its stacking at once; equipartition
gives this term $\frac12\kB T$ on average whatever $N$
(\cref{sec:sm-equipartition}). The typical slide, $\sqrt{\kB T/N\kappa}$,
therefore shrinks as the sheet grows, and the barrier between one
stacking and the next, $N$ times that for one atom, grows. A sheet of 32
atoms wanders much further than the sheets of a crystal, and a larger
cell is the test of how much of what follows depends on it
(\cref{sec:cv-size}).

\subsection*{Counting hops}

The runs start from the relaxed cell with velocities drawn at the target,
and half the kinetic energy passes at once into the potential energy, as
in the crystal of \cref{sec:sm-trajectory}: the first picosecond averages
200, 301 and \SI{402}{\kelvin}. The test of \cref{sec:cv-equilibration}
on the kinetic temperature finds the start at 2.5, 2.5 and
\SI{7.6}{\pico\second}, after which the runs average 298, 447 and
\SI{592}{\kelvin}, and hops are counted from there, over 17.5, 17.5 and
\SI{12.4}{\pico\second}.

A visit counts only if it lasts a least stay, the same device as the
least lifetime of a bond (\cref{sec:ob-bonds}); a shorter visit is joined
to the one before it. With both families of sites the count barely
depends on the least stay up to \SI{50}{\femto\second}: 14, 14, 12 and
12 hops at \SI{300}{\kelvin}, 45, 43, 43 and 42 at 450 and 44, 43, 42
and 42 at \SI{600}{\kelvin}, for least stays of 10, 20, 30 and
\SI{50}{\femto\second} (\cref{fig:ip-hops}b). Beyond, the least stay
begins to discard visits that were real, most at \SI{600}{\kelvin}, where
a visit lasts \SI{0.30}{\pico\second} on average. With the sites of the
sheet below alone the count keeps falling at \SI{300}{\kelvin}, from 28
at a least stay of \SI{10}{\femto\second} to 8 at
\SI{200}{\femto\second}: a position over a site of the other family is
equally far from three sites of the first, and the jitter shares it out
among them, each change a hop. A count that falls steadily as the least
stay grows is the sign of a wrong definition of a site; a plateau is the
sign of a right one, and the least stay is chosen on it, here
\SI{50}{\femto\second}. The hops then join the two families, all of them
at \SI{300}{\kelvin}, 98\% at 450 and 95\% at \SI{600}{\kelvin}, as hops
between neighbours on the honeycomb must; the few others join two sites
of one family, the visit between them having been shorter than the least
stay.

Bonds would not have served. The bond tracker of \cref{sec:ob-bonds},
with its thresholds of 2.45 and \SI{2.86}{\angstrom} for lithium and
carbon, finds lithium bonded to ten carbon atoms in half the frames at
\SI{300}{\kelvin} and to nine in 45\%: the six of the hexagon on one
side, \SI{2.30}{\angstrom} away at \SI{0}{\kelvin}, the atom opposite,
\SI{2.08}{\angstrom}, and that atom's three neighbours,
\SI{2.45}{\angstrom}, on the threshold at which a bond forms. Those three
flicker in and out with the vibrations, each flicker an event of the
tracker, while a hop only exchanges which sheet holds the six and which
the four. The sites count the hops; the bonds would count the
vibrations.

\subsection*{Rates}

With a least stay of \SI{50}{\femto\second} the runs hold 12, 42 and 42
hops. A rate from $n$ events is known to $1/\sqrt n$
(\cref{sec:ob-reactions}), and we call a rate resolved at a fifth, which
needs 25 of them: at \SI{450}{\kelvin} the rate is $(2.40 \pm
0.37)$\,\si{\per\pico\second} and at \SI{600}{\kelvin} $(3.38 \pm
0.52)$\,\si{\per\pico\second}, while at \SI{300}{\kelvin}, $(0.69 \pm
0.20)$\,\si{\per\pico\second} from 12 hops, it is unresolved, and a run of
\SI{36}{\pico\second} in all at that rate would resolve it. The error
$\sqrt n$ assumes hops independent in time. Counted in ten equal windows
of each run, their variance is 1.3, 1.3 and 1.8 times their mean, against
$1 \pm 0.47$ for independent events (\cref{ex:cv-error-error}), within
twice that scatter, so the Poisson errors stand. Over the whole runs,
warming included, the ratios are 1.9, 2.2 and 2.3, beyond twice that
scatter at 450 and \SI{600}{\kelvin}: while a run warms its hops come in
bursts, one more reason to discard the start.

The logarithm of the rate against $1/\kB T$, each point weighted by $\sqrt
n$ as in \cref{sec:ob-reactions}, has the slope $-(80 \pm
17)$\,\si{\milli\electronvolt} (\cref{fig:ip-hops}c): the rate rises 4.9
times from 300 to \SI{600}{\kelvin}, where a barrier of
\SI{73}{\milli\electronvolt} would raise it $\mathrm{e}^{0.073\times(38.68
- 19.34)} = 4.1$ times, with $1/\kB T$ in \si{\per\electronvolt}. The
apparent activation energy lies above the bound from the relaxed path by
less than half its error, which at a fifth of its value is all that three
points, one of them unresolved, can give; a barrier from rates to a few
\si{\milli\electronvolt} needs more temperatures and many more hops at
each.

\begin{figure}[htb]
\centering
\includegraphics{ch17_practice/hops}
\caption{Lithium in dilute graphite, one lithium among 64 carbon atoms,
with MACE-MP-0 and D3. (a) The energy along the straight line to the
nearest site of the other family (teal) and to the next site of the
same family across a C--C bond (ochre), every atom relaxed but lithium
and one carbon of each sheet, whose positions in the plane are held.
(b) Hops counted after each run's start against the least stay a visit
needs to count, with the sites of both families (solid) and of the sheet
below alone (dashed), and the 25 hops that resolve a rate to a fifth
(dotted). (c) The rate of hops for a least stay of
\SI{50}{\femto\second}, with its Poisson error (open: unresolved),
against $1/\kB T$, with the line $\propto\mathrm{e}^{-E_{\mathrm a}/\kB
T}$ through the point at \SI{450}{\kelvin}, for $E_{\mathrm a}$ the
highest point of the teal path (dotted).}
\label{fig:ip-hops}
\end{figure}

\subsection*{Diffusion}

If each hop were independent of the last, the rate would give the
diffusion coefficient, $D = k_{\mathrm h}\ell^2/4$ with $\ell =
\SI{1.4231}{\angstrom}$ (\cref{ex:ip-honeycomb}): with $k_{\mathrm h} =
\num{2.40e-3}$ and \num{3.38e-3} per femtosecond, $D =
\num{2.40e-3}\times1.4231^2/4 =
\SI{12.2e-4}{\angstrom\squared\per\femto\second}$ at \SI{450}{\kelvin} and
\SI{17.1e-4}{\angstrom\squared\per\femto\second} at \SI{600}{\kelvin}.
Lithium's own mean squared displacement after the start, read in the
frame midway between the two sheets and fitted from 0.5 to
\SI{2}{\pico\second} in four blocks, gives $(5.8 \pm 2.2)$, $(22.4 \pm
16.0)$ and $(36.1 \pm 12.6)\times10^{-4}$\,\si{\angstrom\squared\per
\femto\second}, fractional errors of 0.38, 0.71 and 0.35, which miss a
target of 5\% by factors of seven to fourteen. Two
departures from the random walk are measured. The hops are not
independent: a hop returns to the site just left in $0.15 \pm 0.06$ of
cases at both 450 and \SI{600}{\kelvin}, errors $\sqrt{p(1-p)/m}$ for a
fraction $p$ of $m$ pairs of successive hops, three errors below the
$1/3$ of a random walk on the honeycomb, where the site just left is one
of three equally likely; at \SI{300}{\kelvin}, $0.27 \pm 0.13$, the
difference is within its error. Lithium therefore tends to carry on the
way it was going, and the correlation factor of \cref{sec:ob-layers}
exceeds 1, as it did on the model surface. And lithium in a site is
carried by its sheet, whose slides are as long as its hops. One lithium
atom for \SI{20}{\pico\second} gives a rate to 15\% and no diffusion
coefficient: the error of $D$ falls as one over the square root of the
run's length, so 5\% at \SI{450}{\kelvin} needs $(0.71/0.05)^2 = 203$
times the run, or as many lithium atoms in a larger cell, which would
also hold the sheets still.

The runs end with \SI{5}{\pico\second} at fixed energy from the last
frame at \SI{600}{\kelvin}, and another from LiC$_6$ at its mean cell,
each keeping 501 frames with their energies and forces. These are the
trajectories on which planned Chapter~28 trains a learned propagator.

\begin{keyidea}
A site is defined in the frame of the atoms that make it, and a hop by a
least stay chosen where the count is flat; a count that keeps falling as
the least stay grows means the sites are wrong. Hops are counted after
the start that the test of Chapter~15 finds. A rate is reported with its
Poisson error once 25 hops are seen, after a check that they are
independent, otherwise as unresolved with the run that would resolve it.
One atom gives a rate long before it gives a diffusion coefficient.
\end{keyidea}

\notebook{17}{count lithium's hops with both families of sites and with
one, and watch the count against the least stay}

\begin{selfcheck}
At \SI{450}{\kelvin}, how long a run of one lithium atom would resolve the
rate to a tenth?
\end{selfcheck}
